\begin{pastexam}[label={px:nccu-114-dsa-14}]{政大資科所 114 學年度 第 14 題}
（單選；每題 4 分）

(4\%) Which of the following statements is \textbf{false}?
\begin{enumerate}[label=\Alph*.]
\item Every problem which can be solved in polynomial time by deterministic algorithms must be a NP
problem.
\item A problem is NP-hard if it is as least as hard as all NP problems.
\item If any NP-Complete problem can be solved in polynomial time, then all NP-hard problems can be
solved in polynomial time.
\item Every NP-Complete problem must be NP.
\end{enumerate}
\end{pastexam}

\begin{pastsoln}
答案是 C。 B 的「as least」是「at least」的筆誤。

逐項檢查：
\begin{itemize}
\item A.~對。 $\mathrm{P} \subseteq \mathrm{NP}$（資結書定理~28.1）：能在多項式時間解，驗證時直接
重新解一遍、不看證據就好，所以也在 NP 裡。
\item B.~對。 這是 NP 難的直觀說法：NP 裡的每個問題都能在多項式時間歸約到它。
\item C.~錯，這是答案。 某個 NP 完備問題有多項式時間解，可以推出 $\mathrm{P} = \mathrm{NP}$
（資結書定理~28.5），也就是所有 \textbf{NP} 裡的問題（包括全部 NP 完備問題）都有多項式時間解。 但
NP 難的問題不一定在 NP 裡，可以難得多。 例子是停機問題：它是 NP 難的（SAT 可以歸約到它：給一個
公式，造一個「窮舉所有賦值，找到滿足的就停、找不到就無窮迴圈」的程式），但停機問題根本不可判定，
任何演算法都解不了，更不會有多項式時間解。
\item D.~對。 NP 完備的定義就是「在 NP 裡，而且是 NP 難」。
\end{itemize}
\end{pastsoln}
